\origpage{116}{ROBERTSON AND SEYMOUR}
\noindent publishing them separately because we feel that they are of some independent interest.

It is important for our applications to deal with a slightly more general problem. (Conveniently, the generalization makes the problem essentially no harder.) A \textit{division} is a finite set of non-empty pairwise disjoint sets. If $\Delta$ is a division, we abbreviate $\bigcup (\delta \in \Delta)$ by $\bigcup \Delta$. If $G$ is a graph and $\Delta$ is a division with $\bigcup \Delta \subseteq V(G)$, a \textit{realization} of $\Delta$ in $G$ is a subgraph of $G$ which is a forest with $|\Delta|$ components $H(\delta)$ ($\delta \in \Delta$) such that $\delta \subseteq V(H(\delta))$ ($\delta \in \Delta$). We say that $\Delta$ is \textit{feasible} in $G$ if there is a realization of $\Delta$ in $G$. The problem DISJOINT CONNECTING SUBGRAPHS is, given a graph $G$ and a division $\Delta$, with $\bigcup \Delta$ a subset of $V(G)$, to decide if $\Delta$ is feasible in $G$. If each member of $\Delta$ has cardinality two this is clearly equivalent to DISJOINT CONNECTING PATHS. We shall give a polynomial algorithm for this problem when $G$ is planar and there are two regions of $G$ which together are incident with every vertex in $\bigcup \Delta$. Essentially equivalent to this and notationally more convenient is to constrain $G$ to be drawn on a cylinder and $\bigcup \Delta$ to be a subset of the boundary of the cylinder.

Let us fix some terminology. A \textit{surface} is a compact 2-manifold with boundary. In this paper we are only concerned with three surfaces, the disc, the M\"obius band, and the cylinder, which is homeomorphic to
\begin{equation*}
\{(x, y) : 1 \le x^2 + y^2 \le 4\}.
\end{equation*}
We generally denote a surface by $\Sigma$, and its boundary by $bd(\Sigma)$. Any component of $bd(\Sigma)$ is called a \textit{cuff} of $\Sigma$. By a \textit{drawing} of a graph $G$ on $\Sigma$ we shall always mean, as usual, a drawing without crossings. It will often be convenient to identify $G$ with a drawing of $G$, and we do so when we expect it will cause no confusion. In particular, if $G$ is drawn on $\Sigma$ we sometimes regard $G$ as a subset of $\Sigma$. A \textit{region} of $G$ is a connected component of $\Sigma - G$. In any drawing, every point of $\Sigma$ is either a vertex of $G$, or lies in an edge, or lies in a region, and only one of these is true. A graph $G$ is \textit{properly drawn} on $\Sigma$ if no edge of $G$ meets $bd(\Sigma)$. (Thus, $bd(\Sigma)$ only contains vertices of $G$ and points which are in regions.) Graphs may have loops or multiple edges. The \textit{valency} of a vertex is the number of edges incident with it; loops are counted twice. We denote the closure of $X$ by $\bar{X}$.

\sectionhead{2. \textsc{Divisions} \textsc{On} A \textsc{Circle}}

In the next section we shall give an algorithm and a structural characterization for our problem when $G$ is drawn on a disc $\Sigma$ and the division $\Delta$ has $\bigcup \Delta \subseteq bd(\Sigma)$. But first we develop some theory of divisions on a circle. (Our method for this disc case is rather more tortuous then necessary; a
